% A LaTeX (non-official) template for ISAE projects reports
% Copyright (C) 2014 Damien Roque
% Version: 0.2
% Author: Damien Roque <damien.roque_AT_isae.fr>

\documentclass[a4paper,12pt,twoside]{report}
\usepackage[utf8]{inputenc}
\usepackage[T1]{fontenc}
\usepackage[english]{babel}
\usepackage{roboto}
\usepackage{a4wide}
\usepackage{graphicx}
\graphicspath{images/}
\usepackage{subfig}
\usepackage{tikz}
\usetikzlibrary{shapes,arrows}
\usepackage{pgfplots}
\pgfplotsset{compat=newest}
\pgfplotsset{plot coordinates/math parser=false}
\newlength\figureheight
\newlength\figurewidth
\pgfkeys{/pgf/number format/.cd,
set decimal separator={,\!},
1000 sep={\,},
}
\usepackage{ifthen}
\usepackage{ifpdf}
\ifpdf
\usepackage[pdftex]{hyperref}
\else
\usepackage{hyperref}
\fi
\usepackage{color}
\hypersetup{%
colorlinks=true,
linkcolor=black,
citecolor=black,
urlcolor=black}

\renewcommand{\baselinestretch}{1.05}
\usepackage{fancyhdr}
\pagestyle{fancy}
\fancyfoot{}
\fancyhead[LE,RO]{\bfseries\thepage}
\fancyhead[RE]{\bfseries\nouppercase{\leftmark}}
\fancyhead[LO]{\bfseries\nouppercase{\rightmark}}
\setlength{\headheight}{15pt}

\let\headruleORIG\headrule
\renewcommand{\headrule}{\color{black} \headruleORIG}
\renewcommand{\headrulewidth}{1.0pt}
\usepackage{colortbl}
\arrayrulecolor{black}

\usepackage{enumitem}
\setlist{nosep}


\fancypagestyle{plain}{
  \fancyhead{}
  \fancyfoot[C]{\thepage}
  \renewcommand{\headrulewidth}{0pt}
}

\makeatletter
\def\@textbottom{\vskip \z@ \@plus 1pt}
\let\@texttop\relax
\makeatother

\makeatletter
\def\cleardoublepage{\clearpage\if@twoside \ifodd\c@page\else%
  \hbox{}%
  \thispagestyle{empty}%
  \newpage%
  \if@twocolumn\hbox{}\newpage\fi\fi\fi}
\makeatother

\usepackage{amsthm}
\usepackage{amssymb,amsmath,bbm}
\usepackage{array}
\usepackage{bm}
\usepackage{multirow}
\usepackage[footnote]{acronym}

\newcommand*{\SET}[1]  {\ensuremath{\mathbf{#1}}}
\newcommand*{\VEC}[1]  {\ensuremath{\boldsymbol{#1}}}
\newcommand*{\FAM}[1]  {\ensuremath{\boldsymbol{#1}}}
\newcommand*{\MAT}[1]  {\ensuremath{\boldsymbol{#1}}}
\newcommand*{\OP}[1]  {\ensuremath{\mathrm{#1}}}
\newcommand*{\NORM}[1]  {\ensuremath{\left\|#1\right\|}}
\newcommand*{\DPR}[2]  {\ensuremath{\left \langle #1,#2 \right \rangle}}
\newcommand*{\calbf}[1]  {\ensuremath{\boldsymbol{\mathcal{#1}}}}
\newcommand*{\shift}[1]  {\ensuremath{\boldsymbol{#1}}}

\newcommand{\eqdef}{\stackrel{\mathrm{def}}{=}}
\newcommand{\argmax}{\operatornamewithlimits{argmax}}
\newcommand{\argmin}{\operatornamewithlimits{argmin}}
\newcommand{\ud}{\, \mathrm{d}}
\newcommand{\vect}{\text{Vect}}
\newcommand{\sinc}{\ensuremath{\mathrm{sinc}}}
\newcommand{\esp}{\ensuremath{\mathbb{E}}}
\newcommand{\hilbert}{\ensuremath{\mathcal{H}}}
\newcommand{\fourier}{\ensuremath{\mathcal{F}}}
\newcommand{\sgn}{\text{sgn}}
\newcommand{\intTT}{\int_{-T}^{T}}
\newcommand{\intT}{\int_{-\frac{T}{2}}^{\frac{T}{2}}}
\newcommand{\intinf}{\int_{-\infty}^{+\infty}}
\newcommand{\Sh}{\ensuremath{\boldsymbol{S}}}
\newcommand{\C}{\SET{C}}
\newcommand{\R}{\SET{R}}
\newcommand{\Z}{\SET{Z}}
\newcommand{\N}{\SET{N}}
\newcommand{\K}{\SET{K}}
\newcommand{\reel}{\mathcal{R}}
\newcommand{\imag}{\mathcal{I}}
\newcommand{\cmnr}{c_{m,n}^\reel}
\newcommand{\cmni}{c_{m,n}^\imag}
\newcommand{\cnr}{c_{n}^\reel}
\newcommand{\cni}{c_{n}^\imag}
\newcommand{\tproto}{g}
\newcommand{\rproto}{\check{g}}
\newcommand{\LR}{\mathcal{L}_2(\SET{R})}
\newcommand{\LZ}{\ell_2(\SET{Z})}
\newcommand{\LZI}[1]{\ell_2(\SET{#1})}
\newcommand{\LZZ}{\ell_2(\SET{Z}^2)}
\newcommand{\diag}{\operatorname{diag}}
\newcommand{\noise}{z}
\newcommand{\Noise}{Z}
\newcommand{\filtnoise}{\zeta}
\newcommand{\tp}{g}
\newcommand{\rp}{\check{g}}
\newcommand{\TP}{G}
\newcommand{\RP}{\check{G}}
\newcommand{\dmin}{d_{\mathrm{min}}}
\newcommand{\Dmin}{D_{\mathrm{min}}}
\newcommand{\Image}{\ensuremath{\text{Im}}}
\newcommand{\Span}{\ensuremath{\text{Span}}}

\usepackage[style=numeric]{biblatex}
\usepackage{csquotes}
\addbibresource{biblio.bib}

\newtheoremstyle{break}
  {11pt}{11pt}%
  {\itshape}{}%
  {\bfseries}{}%
  {\newline}{}%
\theoremstyle{break}

%\theoremstyle{definition}
\newtheorem{definition}{Définition}[chapter]

%\theoremstyle{definition}
\newtheorem{theoreme}{Théorème}[chapter]

%\theoremstyle{remark}
\newtheorem{remarque}{Remarque}[chapter]

%\theoremstyle{plain}
\newtheorem{propriete}{Propriété}[chapter]
\newtheorem{exemple}{Exemple}[chapter]

\parskip=5pt
%\sloppy


\usepackage{graphicx} % This lets you include figures
\usepackage{hyperref} % This lets you make links to web locations
\graphicspath{ {./images/} }

\usepackage[rightcaption]{sidecap}
\usepackage{wrapfig}

\usepackage{float}

\usepackage{imakeidx}

\usepackage{listings}

\usepackage{pgf-umlsd}

\usepackage{amssymb}

\makeindex




\begin{document}

%%%%%%%%%%%%%%%%%%
%%% First page %%%
%%%%%%%%%%%%%%%%%%

\begin{titlepage}
\begin{center}

\includegraphics[width=0.65\textwidth]{IRT.jpg}\\[1cm]



% Title
\rule{\linewidth}{0.5mm} \\[0.4cm]
{ \huge \bfseries Project 074 : Simulation to Real\\[0.4cm] }

{\large Study on Gym-Duckietown environment, supervised by David Bertoin }\\[0.5cm]

\rule{\linewidth}{0.5mm} \\[1.5cm]

% Author and supervisor
\noindent
\begin{minipage}{0.29\textwidth}
    \begin{center} \large
    \emph{Auteurs :}\\
        Olivier SUTTER\\
        Mathieu VERM\\
        Michael ROMAGNE\\
        Alexandre MARTIN\\
        Pablo MIRALLES\\
        Vincent COYETTE\\
    \end{center}
\end{minipage}%

\smallbreak
\vfill


\smallbreak
\includegraphics[width=0.26\textwidth]{logo-isae-supaero}\\[1cm]

% Bottom of the page
{\large \today}
\end{center}

\end{titlepage}


\tableofcontents

\clearpage
\newpage


\chapter*{Introduction}

Reinforcement Rearning is an area of machine learning. 
The goal is for an agent to learn an optimal behaviour and adapt its comportement, based on experience, in a given environment. 
Reinforcement Learning has been very successful in the last few years, notably for its great performance in games (e.g. AlphaGo\footnote{Silver, D., Schrittwieser, J., Simonyan, K. et al. Mastering the game of Go without human knowledge. Nature 550, 354–359 (2017)}). 
However, it suffers from a lack of industrial applications.

One of the reason for this is that an agent needs a lot of experience in order to learn a reasonable behaviour. 
Those simulations must be run in a simulator for safety, speed and financial reasons. 
Moreover, the transfer of behaviour from the simulated agent to the real world (for example an autonomous vehicle) is far from being trivial.
Even for high quality simulators, there is a shift in both states and transitions spaces between the simulator and the real world.

The purpose of this project is to study different methods to robustly train an autonomous car in a simulator, so that the real-world agent behaves as suited.


DuckieTown\footnote{See https://www.duckietown.org/} environment will be used during all the project.
It provides a simulator (Gym-Duckietown), a real-world agent and a real-world environment. 


\chapter{Duckietown environment}

A reinforcement learning agent usually acts in an environment described by a Markov Decision Process (MDP).
A MDP is basically a 4-tuple $(S, A, P_a, R_a)$.
Let's define each of these components for the duckietown environment.

\section{State}
The state S is the state of the environment from the agent point of view.
In our case, the robot contains only one sensor, the camera.
Thus, the states of this MDP are the simulated images in the simulator or the camera output in the real world.
The camera has a resolution of 640x480, so the state space is $S = \{0, 255\}^{640x480x3}$ (the 3 stands for the 3 RGB colors).

\section{Actions}
The actions are composed by the velocities of each of the two wheels.
By default, the action space is continuous, $A = [-1, 1]^2$, $1$ being full forward velocity for a wheel, $-1$ full backward.

\textbf{Caution :}
Be careful here as the definition of the actions in the README of the duckietown original repository is describing actions differently.
It also defines a two-elements tuple, but the first digit represents the forward velocity while the second one represents the steering angle.
This definition is more natural to manually control the robot.
A wrapper (DuckietownEnv) exists to switch from the wheels velocity control to the forward velocity and streering angle control.

However, this wrapper is not used in this project and the default implementation is kept.
The original code contains some inconsistencies without the wrapper, which will be described later.

\section{Transitions}
The transitions of a MDP is the distribution $P(s'| s, a)$, i.e. the probability of ending up in the state $s'$ if the agent is in the state s and take the action a.
This is typically a point where the transfer from the simulator to the robot may be harmful.

Indeed, the transitions in the simulator can be considered as "perfect", as the position of the robot in state $s'$ is computed from the position in state $s$ using the velocities stated by action $a$.
This is not fully deterministic, as some objects (such as duckie pedestrians) may have a stochastic behaviour, but the updated position is a deterministic function of the previous position and current action.

However, in the real world, the transitions between the positions are not perfect.
Given a position and an action, the next position is not deterministic.
It can slightly vary because of mechanical links not being perfect, because of unusual road surface, etc.

This shift in the transition distribution is one of the problem we will have to address during the simulation to reality transfer.

\section{Reward}
The choice of the reward function is primordial to properly train the agent. The reward may depend on :

\begin{itemize}
    \item The position of the agent. The position may be used to compute the distance from a target position, and the distance from the middle of the lane.
    \item The speed of the agent. The faster the better.
    \item Collisions. The agent may get a negative reward in case of collision, and possibly if it gets too  close to another object.
\end{itemize}

The default reward in duckietown is as follows:

\begin{align}
R(t) = 40 * C_{p} - 10*dist + alignment*speed \\
C_p = \sum_{Obj} ((pos_{agent} - pos_{object}) - SR_{agent} - SR_{object}) \\
\end{align}
\begin{itemize}
    \item $C_p$ is the collision penalty for being dangerously close to other objects. It is a proxy for area overlap. Notably, one can collide with several objects at once (with additive effects) and one can collide with respect to the reward function (which uses the safety radius) without the episode restarting (which depends on the collision radius). Note that collision penalty is smaller than 0 whenever there is a collision, therefore this is actually a penalty despite the positive sign.
    \item dist is the distance between the agent's center and the closest point on the line defining the lane's center
    \item alignment is the dot product between direction and the normalized tangent of the road.
    \item speed is the agent's speed
\end{itemize}

This reward function has presented several issues:

\begin{enumerate}
    \item The penalties (collision and distance) are so strong that the agent usually much prefers to go around in circles at max speed at the center of the lane than to actually follow it.
\end{enumerate}


The design of the reward function will condition the performance of our agent in the simulator.
This reward will be used to train a policy, i.e. a mapping from a state $s$ to an action $a$ to take.
This policy will need to be adapted before transfer to the real world.
However, the reward does not need to be adapted.

\section{Policy}
A policy $\Pi$ associates to each state a distribution of probability over the actions, which is used to select the action to take.
The training of the agent aims at finding a policy selecting the action which maximize the reward from each state.
Such policy is named optimal.
A deterministic optimal policy is supposed to exist, and thus the policy will be searched as a mapping from the states space to the action space.


\chapter{Gym-DuckieTown Simulator}

From the README.md in the github repository \footnote{$https://github.com/duckietown/gym-duckietown$}, here is an introduction to Gym-Duckietown. 

 \begin{quotation}
Gym-Duckietown is a simulator for the Duckietown Universe, written in pure Python/OpenGL (Pyglet). It puts the agent, a Duckiebot, inside of an instance of a Duckietown: a loop of roads with turns, intersections, obstacles, Duckie pedestrians, and other Duckiebots. It can be a pretty hectic place!
 \end{quotation}

This project has a dedicated github repository\footnote{$https://github.com/vcoyette/gym-duckietown$}, which is a fork from the Gym-Duckietown original repo.

\section{Installation}
The installation is pretty straight-forward from the repository. Use the following commands :

\begin{lstlisting}[language=bash]
    git clone https://github.com/vcoyette/gym-duckietown
    cd gym-duckietown
    conda env create -f environment.yaml
\end{lstlisting}

The installation has been tested on Windows, Linux and MacOS. Some problems may be encountered for the installation of certain packages. They can be resolved with package-specific installation instructions. 
For example, the installation of pyglet may raise an issue. It can be resolved by installing it from the pyglet github repository. 

To use the simulator, the environment must be activated :
 
\begin{lstlisting}[language=bash]
    conda activate gym-duckietown
\end{lstlisting}

And the root folder of the project must be added to the PYTHONPATH environment variable.
On linux :
\begin{lstlisting}[language=bash]
    export PYTHONPATH="${PYTHONPATH}:`pwd`"
\end{lstlisting}

On Windows, environment variable can be accessed in the parameters,  advanced parameters. You can then append the path of your project folder to the PYTHONPATH variable if it exists, or create it otherwise.


\section{Manual Control}
A UI application can be launched to manually control the robot. Actions can be sent from the keyboard, and images from the DuckieBot camera are displayed. Here is a simple command to launch the application :

\begin{lstlisting}{laguage=bash}
$ ./manual_control.py --env-name Duckietown-udem1-v0
\end{lstlisting}

\subsection{Parameters}
Here is a list of parameters which can be used. 

\begin{itemize}
	\item env-name: the name of the environment to execute (TODO anchor to env description)
	\item map-name: the name of the map to be used
	\item distortion: boolean, add distorsion
	\item draw-curve: boolean, draw the lane-following curve
	\item draw-bbox: boolean, draw the collision bounding box
	\item domain-rand: boolean, use domain randomization
	\item frame-skip: number of frames to skip (default 1)
	\item seed: seed
\end{itemize}

\subsection{Keyboard}
Here is a list of keys which can be used during the simulation.

\begin{itemize}
	\item Escape : exit simulation
	\item Backspace : restart simulation
	\item Directional keys: go forward, backward or turn
	\item Shift : boost speed
\end{itemize}

\subsection{Logs}
The github repository contains a branch "experiment". This branch is to be used to do any experiment on the simulator. The main difference from a classical "develop" branch is that it does not aim at being merged into the master branch. 

In this branch, we designed a logger for the manual control application. At each time step, this logger will log the current position of the agent, the current speed, the current distance from the lane, the action took. The goal is to manually control the agent so that it behaves as we would expect. We would then be able to check the logs, and use them to design a reward function.

To enable this option, pass the --output option to manual\_control.py. You can optionally specify another option `--filename example.csv` to specify the name of the output file, which would be in data/example.csv. If no file-name is specified, the logs will be stored in data/manual\_controli.csv, where i is the first number for which this path is free.

When the episode is done, the manual control must be exited by pressing the S key to save the output.

The learning is performed using the train\_reinforcement.py script. In order to execute it, change into the directory learning and run : 

\begin{lstlisting}{laguage=bash}
    python -m reinforcement.pytorch.train_reinforcement
\end{lstlisting}

If an error is returned stating that a module doesn't exist, check if PYTHONPATH variable contains the repository base folder.

\begin{lstlisting}{laguage=bash}
    echo $PYTHONPATH
\end{lstlisting}

If it doesn't, get back to your to the root directory of the project and run: 

\begin{lstlisting}{laguage=bash}
    export PYTHONPATH="${PYTHONPATH}:`pwd`"
\end{lstlisting}

The default algorithm implemented is DDPG. It can be tweaked by passing some hyperparameters when executing train\_reinforcement. A list of the parameters can be accessed by running :

\begin{lstlisting}{language=bash}
    python -m reinforcement.pytorch.train_reinforcement --h
\end{lstlisting}

\section{Create Maps}
You can very easily create a new \textbf{Duckietown} environment with a text editor. A Duckietown's map is a \textit{.yaml} file, so you have to save your new map in the folder \textbf{maps} as "\textit{my\_new\_map.yaml}" (the path should looks like this one : \textit{./gym-duckietown-master/gym\_duckietown/maps}).
\newline

If you want to see your map, you can use the following line in the terminal : \newline

\begin{lstlisting}{language=bash}
./manual_control.py --env-name Duckietown-udem1-v0 
--map-name my\_new\_map}
\end{lstlisting}

\noindent Your robot will be manually controlled in your map.\newline


\noindent A grassroots level of your \textit{.yaml} should looks like this :\newline
\_\_\_\_\_\_\_\_\_\_\_\_\_\_\_\_\_\_\_\_\_\_\_\_\_\_\_\_\_\_\_\_\_\_\_\_\_\_\_\_\_\_\_\_\_\_\_\_\_\_\_\_\_\_\_\_\_\_\_\_\_\_\_\_\_\_\_\_\_\_\_\_\_\_\_\_\_\_\_\_\_\_\_\_\_\_\_\_\_\_\_\_\_\_\_\_

\noindent tiles:\newline

\noindent- [floor, floor, floor        , grass        , grass        , grass, floor, floor]\newline
\noindent- [floor, floor       , grass        , grass        , straight/S   , grass, floor, floor]\newline
\noindent- [floor, grass       , grass        , curve\_left/W , curve\_right/S, grass, floor, floor]\newline
\noindent- [grass, grass       , curve\_left/W , curve\_right/S, grass        , grass, floor, floor]\newline
\noindent- [grass, curve\_left/W, curve\_right/S, grass        , grass        , floor, floor, floor]\newline
\noindent- [grass, curve\_left/S, curve\_right/E, grass        , grass        , floor, floor, floor]\newline
\noindent- [grass, grass       , curve\_left/S , curve\_right/E, grass        , floor, floor, floor]\newline
\noindent- [floor, grass       , curve\_left/W , curve\_right/S, grass        , floor, floor, floor]\newline
\noindent- [floor, grass       , straight/S   , grass        , grass        , floor, floor, floor]\newline
\noindent- [floor, grass       , grass        , grass        , floor        , floor, floor, floor]\newline
\noindent- [floor, floor       , floor        , floor        , floor        , floor, floor, floor]\newline

\noindent objects:\newline

\noindent- kind: house\newline
  pos: [4.5, 9.1]\newline
  rotate: 90\newline
  height: 0.5\newline

\noindent- kind: tree\newline
  pos: [1, 1]\newline
  rotate: 0\newline
  height: 0.5\newline

\noindent- kind: tree\newline
  pos: [2, 8.5]\newline
  rotate: 90\newline
  height: 0.3\newline

\noindent tile\_size: 0.585\newline
\_\_\_\_\_\_\_\_\_\_\_\_\_\_\_\_\_\_\_\_\_\_\_\_\_\_\_\_\_\_\_\_\_\_\_\_\_\_\_\_\_\_\_\_\_\_\_\_\_\_\_\_\_\_\_\_\_\_\_\_\_\_\_\_\_\_\_\_\_\_\_\_\_\_\_\_\_\_\_\_\_\_\_\_\_\_\_\_\_\_\_\_\_\_\_\_\newline

For each line, the number of tiles has to remain the same. The \textbf{tile\_size} and the \textbf{height} of every object can change, but insofar as your goal is to transpose the simulation to the real world, you must not change their values. If so, you should have \textbf{tile\_size}$=0.585$ (for the conventional heights of the objects will be given below).\newline

You can find your way in the map knowing that going upward is going North, and knowing that when you drive on a road tile, the name of the road tile tells you which direction you were facing when you arrived on it, and by which direction you will leave the tile. Let's give an example : the tile "\textit{curve\_left/W}" means that you arrived on it while you were moving West, and the fact that you'll turn left on it says you are going to leave the tile by the South (it's of course reversible, you can arrive from South facing North, and leave to the East).\newline
You can also know your position $(x,\ y)\in\mathbb{R}^2$ on the map (this is the way you put objects on it). The point $[0,\ 0]$ matches the upper left corner of the upper left tile of the map (so the coordinates start at the very North/West). The coordinates keep going higher as you move toward South/East, increasing of 1 each time you go through one tile.\newline

Here is the list of the tiles you can use to build your map:

\noindent- straight\newline
\noindent- curve\_left\newline
\noindent- curve\_right\newline
\noindent- 3way\_left (3-way intersection)\newline
\noindent- 3way\_right\newline
\noindent- 4way (4-way intersection)\newline
\noindent- asphalt\newline
\noindent- grass\newline
\noindent- floor (office floor)\newline

You can (and should) orientate the roads (the six first tiles on the list) by adding "/N", "/E", "/S", "/W" (this will be oriented according to the rule given above).\newline

The objects can be added as shown in the following example (changing the name of the object). Here is a list of them :\newline

\noindent- barrier (height : 0.08)\newline
\noindent- cone (height : 0.08)\newline
\noindent- duckie (height : as you wish between 0.06 and 0.08)\newline
\noindent- duckiebot (height : 0.12)\newline
\noindent- tree (height : as you wish between 0.1 and 0.9)\newline
\noindent- house (height : 0.5)\newline
\noindent- truck (height : 0.25)\newline
\noindent- bus (height : 0.18)\newline
\noindent- building (height : 0.6)\newline
\noindent- sign\_stop, sign\_T\_intersect, sign\_yield, etc... (height : $0.18$ for the signs, $0.4$ for a traffic light)\newline

There are many other signs, you can check the whole list here :\newline
https://github.com/duckietown/gym-duckietown/blob/master/gym\_duckietown/meshes\newline

It is possible to add the attributes : \newline

\noindent- optional: True or False (makes the object optional)\newline
\noindent- static: True or False (for the Duckiebot for example if you want to see them move)\newline

To go any further about map creation, check the Github of Duckietown on this link : \footnote{$https://github.com/vcoyette/gym-duckietown/tree/documentation$ }\newline

\section{Reward Wrapper}
The reward function is computed by the compute\_reward function defined in the class Simulator (in the file gym\_duckietown/simulator.py). The best way to modify the reward function is to edit the function reward of the DtRewardWrapper class in the file learning/utils/wrappers.py.

This class is a subclass of the gym RewardWrapper class. This class implements the Wrapper/Adapter design pattern. It overrides the step function of the simulator. It is hard to explain, but you can have a look at the following sequence diagram to get an idea of how it works.

The idea is that the reward wrapper replaces the environment in the DDPG implementation. During the training, the DDPG calls the step function of the wrapper. This wrapper instance has an attribute which is the environment. The step function in the wrapper calls the function step of the environment, and captures the returned values (observation, reward, done and info). Then, the wrapper runs its own reward function with the reward returned by the previous call as parameter. The wrapper step function returns the values of observation, done and info which were returned by the environment, and returns its custom reward to replace the reward of the environment.

The wrapper containing an environment attribute, it has access to the same information as the simulator to compute its custom reward. For example, it can use the current position through env.cur\_pos, the angle through env.angle, etc.


\begin{sequencediagram}
    \newthread{t}{DDPG}
    \newinst[1]{w}{DtRewardWrapper}
    \newinst[2]{s}{Simulator}
    \begin{call}{t}{step()}{w}{\shortstack{
        return obs, \\ reward(reward) \\
        done, \\ info }}


        \begin{call}{w}{step()}{s}{\shortstack{
            return obs, \\ reward, \\
            done, \\ info }}
            \postlevel
            \postlevel
            \postlevel
        \end{call}
    \postlevel
    \postlevel
    \postlevel
    \end{call}
\end{sequencediagram}

TODO : add legend to diagram











\section{Domain randomization}

When it comes to transfer knowledge from simulation to reality, a problem we may face is that the images collected from simulation diverge too much from reality. Often, people even retrain their model from scratch when moving to the physical world.\\
One solution to this problem is domain randomization. The idea is to perturb the dynamics or look of the simulator such as colors, textures, horizon ... Thus, we obtain a more variable dataset and we have a better generalization, which is beneficial for transfer on the real robot. This part will deal with how you can do domain randomization in gym-duckietown environments.



\subsection{The randomisation API}
The folder at \url{gym-duckietown/gym_duckietown/randomization} contains the domain randomization API. This API contains all of the pre-packaged methods for randomization within gym-duckietown, which I’ll list here. This folder also contains a readme file detailing the API.

The domain randomization is driven by the Randomizer class, which takes as input a configuration file and outputs (upon call to randomize) a set of settings used by the Simulator class (the core class of gym\-duckietown managing the environment). To activate any domain randomization at all, the simulator class must have domain\_rand \= true passed as a parameter to its constructor

If a randomizable variable is not found in the configuration file, it will be randomized according to the default values found in \url{gym-duckietown/gym_duckietown/randomization/config/default.json}. 3 types of distribution are supported, int, uniform and normal. 4 variables are randomizable by default:
\begin{enumerate}
    \item horz\_mode: The task is made harder by making the horizon more similar to the road. It can take integer values from 0 to 3 where:
    \begin{enumerate}
        \item 0: Sets the skybox to a blue sky, the default
        \item 1: Sets the skybox to a gray wall, intended to be similar to room testing conditions
        \item 2: Sets the skybox to a dark gray box
        \item 3: Sets the skybox to a light gray box
    \end{enumerate}
    \item light\_pos: Makes the simulator more or less iluminated, by changing the position of the single light source.
    \item camera\_noise: Adds noise to the camera position for data augmentation purposes. This noise is applied at each render, giving the screen input a "twitchy" behaviour.
    \item frame\_skip: No info provided. Code inspection shows this is the number of frames to skip per action. Higher frameskip makes the agent able to act less frequently.
\end{enumerate}

The API is sorely lacking in the amount of variables that can be randomized. It provides some flexibility in randomization of the input, but it provides no randomization besides frame\_skip of transitions and/or actions.
\subsection{Non API Randomisation}
Code inspection reveals that more randomisation occurs, contingent on domain-rand being activated. 
\begin{enumerate}
    \item Within Simulator.py: All calls to \_perturb(val,scale) distort the value passed as argument if and only if domain rand is true, otherwise they return it undisturbed. The distortion is a multiplicative \% error drawn uniformly between $1 - scale$ and $1 + scale$, with a default of $scale = 0.1$ (e.g. by default values are randomised between 90 and 110\%). The function is called on:
    \begin{itemize}
        \item horizon\_color, beyond the randomization introduced by horz\_mode. 10\% for blue\_sky and wall\_color sky (modes 1 and 2), 40\% for modes 3 (dark gray) and 4 (clear gray)
        \item glLight function : sets the values of individual light source parameters, making the environment more or less iluminated. It modifies the light position, the ambient and diffusion of light.
        \item ground\_color, 30\%
        \item wheel\_dist, 10\%
        \item cam\_height, 8\%
        \item cam\_angle, 20\%
        \item cam\_fov\_y, 20\%, camera field of view
        side length of the ground/noise triangles generated as distractors (which themselves are generated randomly in the first place? Seems redundant to do it twice), 10\%
        \item tile color, 20\% for each tile
        \item object color, 20\% for each color
        \item CAMERA\_FORWARD\_DIST of gl.glTranslatef on line 1434, 10\%
        \item The actual tile texture loaded is randomised with randint amongst all possible candidates
        \item Some optional objects are invisible, 33\% chance
    \end{itemize}
    \item Within objects.py
    \begin{itemize}
        \item DuckiebotObj (Cars)
        \begin{itemize}
            \item follow\_dist is randomised from 0.3 to a uniform between 0.3 and 0.4
            \item velocity is randomised from 0.1 to a uniform between 0.05 and 0.15
        \end{itemize}
        \item DuckieObj (Pedestrians)
        \begin{itemize}
            \item pedestrian\_wait\_time randomised from 8 to a randint from 3 to 20, takes on new random value on same range when finish crossing street
            \item vel randomised from 0.02 to a normal with avg 0.02 and stdev 0.005, takes on new random value on same range when finish crossing street
        \end{itemize}
        \item TrafficLightObj
        \begin{itemize}
            \item freq is randomised from 5 with randint(4,7)
            \item The lights start randomly from off as either On or Off, 50\% chance each
        \end{itemize}
    \end{itemize}
\end{enumerate}


\subsection{Randomizing inputs}


Randomizing inputs can be accomplished via the API. Other possible sources of randomization are:
\begin{itemize}
    \item Changing hue of key elements: stop signs, road… Might be hard since they’re all texture based. A possible approach: Generating variations on existing textures: Automatically generate X diff textures with some other software from existing textures passed through some filters, then use the in-built texture randomiser
    \item Adding noise to the camera acquisition (not position)? This is similar to randomising the color of the sky and objects themselves so it might not be interesting. However camera noise provides variations during a single episode or each timestep which may imitate real camera noise (or not?)
    \item Approaching photorrealism in some way can be helpful in improving the performance (Johnson-Robertson et al.) May prove unfeasible given the simulator and available time and resources
    \item Addding additional light sources?
    \item Gaussian noise foreground of images and edge blurring through gaussian noise at edges on input images is also helpful (Hinterstoisser et al.) The technique can be extended to other blending techniques (Dwibedi et al.)
\end{itemize}





\chapter{Reinforcement Learning Training}%

\section{Training Scipt}%

The learning is performed using the train\_reinforcement.py script in the learning/reiforcement/pytorch directory.
In order to execute it, change into the directory ./learning and run :

\begin{lstlisting}{laguage=bash}
    python -m reinforcement.pytorch.train_reinforcement
\end{lstlisting}

If an error is returned stating that a module doesn't exist, check if PYTHONPATH variable contains the repository root folder.

\begin{lstlisting}{laguage=bash}
    echo $PYTHONPATH
\end{lstlisting}

If it doesn't, get back to the root directory of the project and run:

\begin{lstlisting}{laguage=bash}
    cd ../
    export PYTHONPATH="${PYTHONPATH}:`pwd`"
\end{lstlisting}

The training can be tweaked by passing some hyperparameters when executing train\_reinforcement.
A list of this parameters can be accessed by running :

\begin{lstlisting}{language=bash}
    python -m reinforcement.pytorch.train_reinforcement --h
\end{lstlisting}

\section{Policy}%

The policy can be tweaked through the "--policy" argument, accepting as a string the name of the policy to use to train the agent.
Currently available values are ddpg\cite{ddpg} and td3\cite{td3}, ddpg being the default.

The policies are defined in their own files (e.g. ddpg.py and td3.py).
The implementation is decoupled from the architecture of actors and critic, which are defined in actor.py and critic.py

Among these paramters, you can notice the "--per" parameter, which can be used to use a Prioritize Experience Replay buffer.
It is for now functional with the ddpg algorithm, \textbf{not with the td3}.



\chapter{Transfer to the Duckiebot}

\section{Assembly and operation of the duckiebot}

For the most part we recommend following the duckiebot operating manual at \url{https://docs.duckietown.org/DT19/opmanual\_duckiebot/out/index.html}.
The guide provides a step by step process on how to succesfully assemble, boot, and operate the duckiebot.

While the guide is helpful, Section D does not provide help on how to troubleshoot many of the issues that can appear during boot. In this section we provide information on how to deal with the problems we faced However we recommend heading to the duckietown slack community over at \url{duckietown.slack.com}, which as is very active as of March 2020 and proved invaluable to us. Their considerable expertise and willingness to help can aid you in troubleshooting your issues.

\subsection{Assembly of the duckiebot}

Note: We strongly recommend flashing the SD card and performing the first boot before mounting the Raspberry Pi on the duckiebot. Connect a monitor and/or a keyboard to the raspberry pi to aid you in this procedure, otherwise your sole source of information on the boot process will be the red and green LED lights, until you are able to connect the bot to the local network and ssh into it. After mounting the raspberry PI, the HDMI and keyboard input ports will be rendered inaccessible.

\subsubsection{Inserting and extracting the SD card after mounting the Raspberry PI}

Mounting will render the SD card port of the Raspberry PI hard to access. We recommend using a set of calipers to access the SD card and insert or extract it.


\subsection{Flashing the SD Card and booting the Raspberry Pi}

daffy is the experimental version. It contains numerous duckietown shell commands available for it exclusivel, that can simplify troubleshooting significantly. dts-19 is the last stable build, although it is comparatively lacking in capabilities.

Here are some of the issues that we've encountered upon setting up the duckiebot and how to solve them

\subsubsection{Booting the Raspberry PI before assembly}

The first boot process of raspberry PI is often faulty. As mentioned at the beginning of the section, to ease troubleshooting, we recommend performing this process BEFORE mounting the Raspberry PI onto the duckiebot to make the procedure easier to follow
\subsubsection{Interpreting the LED lights}
The duckiebot operating manual mentions that the following pattern is to be observed during a successful boot of the duckiebot:

\begin{enumerate}
    \item Stable red LED light and intermittent green LED light
    \item Alternating red and green LED lights
    \item Persistent green LED light - The process succeeded
    \item A persistent red LED light means the process failed somehow
\end{enumerate}

The guide does not explain properly that this process is ONLY to happen during the FIRST boot of the raspberry PI. Subsequent boots will have the red LED light lit permanently whereas the green LED light just marks activity on the SD port. Do not be alarmed by this.

\subsubsection{Failed boot procedure - containers missing}

Very often (100\% of the time so far), the boot procedure will fail silently - several of the required packages will fail to be installed with no apparent error message during the flashing procedure or on the duckiebot's health logs. Should that be the case, we recommend that you try to download these containers individually from docker and set them up yourself. The specifics of this process vary case by case. The only alternative is to reflash the SD card and reboot the Raspberry Pi, a time-consuming procedure with low success probabilities.

ATTACH IMAGE OF ALL CONTAINERS NEEDED FOR REGULAR OPERATIONS

\subsection{Loading an agent}

TO BE COMPLETED

\subsection{Calibrating the camera}

TO BE COMPLETED

\subsection{Calibrating the wheels}

TO BE COMPLETED

\end{document}


\printbibliography

\end{document}



